\begin{table}[H]
\centering
\renewcommand{\arraystretch}{1.3}
\begin{tabular}{L{3.3cm}L{8.7cm}}
\toprule
\textbf{Field} & \textbf{Role} \\
\midrule
\texttt{text}              & Raw message text, natural Vietnamese with common code-switching. \\
\texttt{label}             & \textbf{Ground-truth class} for the classification task: one of \texttt{bank\_impersonation}, \texttt{zalo\_social\_engineering}, \texttt{task\_scam}, \texttt{benign}. Set during generation. \\
\texttt{risk\_tier}        & How dangerous the message is: \texttt{benign}, \texttt{suspicious}, or \texttt{high-risk}. Decided separately from the label, not calculated from it. \\
\texttt{suspicious\_spans} & The exact phrases inside \texttt{text} that justify the label: fake links, phone numbers, pressure wording. \\
\texttt{xai\_explanation}  & A short written reason for the label, which is what teaches the model to explain itself. \\
\texttt{source}            & Which model wrote the message, so every message can be traced to its generator. \\
\texttt{seed\_id}          & Links the message back to the \textit{tinnhiemmang.vn} article it was written from. Messages sharing this value always go into the same split. \\
\bottomrule
\end{tabular}
\caption{The seven fields every message in the dataset carries. Pydantic checks all of them before a message is accepted.}
\label{tab:record-schema}
\end{table}
